%% ============================================================
%%  SECCIÓN 1 — Introducción
%%  Responsable: [Minilux]
%% ============================================================

\section{Introducción}

El estudio de la elasticidad constituye uno de los pilares fundamentales de la mecánica
de materiales. Cuando un cuerpo sólido es sometido a una fuerza externa, experimenta
una deformación cuya naturaleza depende tanto de las propiedades intrínsecas del
material como de la magnitud del esfuerzo aplicado. Si el cuerpo recupera su forma
original al retirar la fuerza, se dice que el comportamiento es \textit{elástico}; en caso
contrario, la deformación es permanente y el material se denomina \textit{inelástico}
\cite{serway}.

La relación entre el esfuerzo aplicado y la deformación resultante fue establecida
empíricamente por Robert Hooke en 1658. Hooke observó que, para deformaciones
pequeñas, la deformación unitaria $\varepsilon$ es proporcional al esfuerzo $\sigma$
ejercido sobre el cuerpo. Esta proporcionalidad, conocida como la \textit{Ley de Hooke},
se expresa como:
\begin{equation}
    \sigma = Y\varepsilon,
\end{equation}
donde la constante de proporcionalidad $Y$ recibe el nombre de \textit{módulo de Young},
en honor al científico Thomas Young quien lo calculó formalmente en 1802. El módulo de
Young caracteriza la rigidez del material y tiene dimensiones de presión (\unit{\pascal})
\cite{halliday}.

El esfuerzo técnico $\sigma_0$ se define como el cociente entre la fuerza aplicada $F$ y
el área de la sección transversal original $S_0$ del cuerpo:
\begin{equation}
    \sigma_0 = \frac{F}{S_0}.
    \label{eq:esfuerzo}
\end{equation}
Asimismo, la deformación unitaria $\varepsilon$ se define como el cambio porcentual de la
longitud respecto a la longitud natural $\ell_0$:
\begin{equation}
    \varepsilon = \frac{\Delta\ell}{\ell_0}.
    \label{eq:deformacion}
\end{equation}
La validez de la Ley de Hooke se restringe a deformaciones pequeñas, dentro del
denominado \textit{límite elástico}. Superado este umbral, la relación entre esfuerzo y
deformación deja de ser lineal y el material puede alcanzar su \textit{punto de ruptura}.

El presente experimento estudia el comportamiento elástico de dos materiales con
propiedades físicas distintas: un resorte metálico y una tira de jebe (liga). Para cada
material se aplican cargas crecientes y se registran las deformaciones producidas, tanto
en la carga como en la descarga, con el fin de determinar experimentalmente la constante
recuperadora del resorte y el módulo de Young a partir de un ajuste lineal por mínimos
cuadrados. La comparación entre los comportamientos de ambos materiales permite además
ilustrar los conceptos de elasticidad, histéresis y límite elástico.